\DocumentMetadata{
  pdfversion=2.0,
  pdfstandard=ua-2,
  lang=en-US,
  tagging=on,
  tagging-setup={math/setup=mathml-SE,tabsorder=structure}
}
\documentclass[11pt,letterpaper]{article}
\usepackage[margin=0.68in,includefoot]{geometry}
\usepackage{newcomputermodern}
\usepackage{amsmath,amscd,mathtools}
\usepackage{tikz,adjustbox}
\providecommand{\square}{\mdlgwhtsquare}
\usepackage[hidelinks]{hyperref}
\hypersetup{
  pdftitle={Probability PhD Exam, May 2008},
  pdfauthor={University of Florida Department of Mathematics},
  pdfsubject={Graduate mathematics examination},
  pdfdisplaydoctitle=true,
  bookmarksopen=true
}
\setcounter{secnumdepth}{0}
\setlength{\parindent}{0pt}
\setlength{\parskip}{0.28em}
\setlength{\emergencystretch}{3em}
\setlength{\leftmargini}{2.15em}
\setlength{\leftmarginii}{2.65em}
\setlength{\leftmarginiii}{3em}
\pagestyle{empty}
\raggedbottom
\allowdisplaybreaks
\NewTaggingSocketPlug{block/list/label}{examproblem}{%
  \tagstructbegin{tag=Lbl}\tagstructend%
  \tagstructbegin{tag=\LBody}%
  \tagstructbegin{tag=\ProblemHeadingTag,title-o={\ProblemHeadingText}}%
  \tagmcbegin{tag=\ProblemHeadingTag,actualtext={\ProblemHeadingText}}%
  #2%
  \tagmcend%
  \tagstructend%
}
\newenvironment{examproblems}[2]{%
  \def\ProblemHeadingTag{#2}%
  \AssignTaggingSocketPlug{block/list/label}{examproblem}%
  \begin{enumerate}%
  \setcounter{enumi}{#1}%
  \renewcommand{\labelenumi}{\textbf{\theenumi.}}%
  \setlength{\topsep}{0.35em}%
  \setlength{\partopsep}{0pt}%
  \setlength{\itemsep}{0.48em}%
  \setlength{\parsep}{0.18em}%
}{\end{enumerate}}
\newenvironment{examsubparts}{%
  \AssignTaggingSocketPlug{block/list/label}{default}%
  \begin{enumerate}%
  \setlength{\topsep}{0.18em}%
  \setlength{\partopsep}{0pt}%
  \setlength{\itemsep}{0.16em}%
  \setlength{\parsep}{0.1em}%
}{\end{enumerate}}
\newenvironment{examinstructions}{%
  \par\begingroup\itshape
}{%
  \par\endgroup\vspace{0.18em}
}
\makeatletter
\renewcommand\subsection{\@startsection{subsection}{2}{0pt}%
  {-1.25ex plus -.25ex minus -.1ex}%
  {0.55ex plus .1ex}%
  {\normalfont\large\bfseries}}
\renewcommand{\@maketitle}{%
  \newpage\null\vskip -1em%
  {\footnotesize\bfseries
    \href{https://www.ufl.edu/}{University of Florida}, \href{https://math.ufl.edu/}{Department of Mathematics}\par}%
  \vskip -0.5em%
  \tagstructbegin{tag=H1,title={Probability PhD Exam, May 2008}}%
  \tagpdfparaOff%
  \tagmcbegin{tag=H1}%
    {\LARGE\bfseries\href{https://gma.math.ufl.edu/exams/probability-phd/}{Probability PhD Exam}, May 2008\par}%
  \tagmcend%
  \tagpdfparaOn%
  \tagstructend%
  \vskip 0.25em%
  \tagmcbegin{artifact}%
    \hrule height 1pt%
  \tagmcend%
  \vskip 1em%
}
\makeatother
\title{Probability PhD Exam, May 2008}
\author{}
\date{}
\begin{document}
\maketitle
\thispagestyle{empty}
\begin{examproblems}{0}{H2}
\def\ProblemHeadingText{Problem 1.}
\item
Let \(\mathcal F\) be a sigma field. Two random variables~\(X\) and~\(Y\) satisfy
\[
\mathbb E[Y^2\mid\mathcal F]=X^2,\qquad \mathbb E[Y\mid\mathcal F]=X.
\]
 Show that \(X=Y\) a.s.
\def\ProblemHeadingText{Problem 2.}
\item
Let~\(X\) and~\(Y\) be two independent random variables with the same probability distribution \(F(x)\). If \((X+Y)/\sqrt{2}\) has the same probability distribution \(F(x)\), then
\[
F(x)=\int_{-\infty}^{x}\frac{1}{\sqrt{2\pi}}e^{-y^2/2}\,\mathrm{d}y.
\]
\def\ProblemHeadingText{Problem 3.}
\item
Let \(f_n(x_1,\ldots,x_n)\) be the probability density function of the random vector \((X_1,\ldots,X_n)\) and \(g_n(x_1,\ldots,x_n)\) be the probability density function of the random vector \((Y_1,\ldots,Y_n)\). We define
\[
Z_n=\frac{g_n(X_1,\ldots,X_n)}{f_n(X_1,\ldots,X_n)},
\]
 if \(f_n(X_1,\ldots,X_n)>0\), otherwise \(Z_n=0\). Show that \(\{Z_n,n=1,2,\ldots\}\) is a supermartingale.
\def\ProblemHeadingText{Problem 4.}
\item
Let \(X_1,X_2,\ldots\) be a sequence of i.i.d random variables with mean~\(0\) and variance \(\sigma^2<\infty\). If \(S_n=X_1+X_2+\cdots+X_n\), then
\[
\lim_{n\to\infty}\mathbb E\left[\frac{|S_n|}{\sqrt n}\right]=\sqrt{\frac{2}{\pi}}\sigma.
\]
\def\ProblemHeadingText{Problem 5.}
\item
Let \(B_t\) be a standard Brownian motion and define
\[
D_n=\max_{n\le t\le n+1}|B_t-B_n|.
\]
 Prove that \(D_n/n\) converges to~\(0\) with probability one.
\def\ProblemHeadingText{Problem 6.}
\item
Consider square integrable random vectors
\[
Y=\bigl(y(1),y(2),\ldots,y(d)\bigr)
\]
 and
\[
Y_n=\bigl(y_n(1),y_n(2),\ldots,y_n(d)\bigr),\qquad n\in\mathbb N,
\]
 which satisfy
\[
\lim_{n\to\infty}\mathbb E\bigl(|y_n(j)-y(j)|^2\bigr)=0\qquad\text{for }j=1,2,\ldots,d.
\]
 Show that the mean of \(Y_n\) converges to the mean of~\(Y\), and the covariance matrix of \(Y_n\) converges to that of~\(Y\).
\def\ProblemHeadingText{Problem 7.}
\item
Suppose that \(\{X_t\}\) and \(\{Y_t\}\) are both standard Brownian motion and are independent. Let \(\mathcal F_t\) denote the~\(\sigma\)-field \(\sigma(X_s,Y_s\mid s\in[0,t])\). Prove that
\[
Z_t=X_t^2Y_t-\int_0^t Y_u\,\mathrm{d}u
\]
 is an \(\mathcal F_t\)-martingale.
\def\ProblemHeadingText{Problem 8.}
\item
Let \(\mathcal F\) be a sub-sigma field of \(\mathcal G\) and~\(Z\) be a nonnegative random variable defined on probability space \((\Omega,\mathcal G,\mathbb P)\) with \(\mathbb E_{\mathbb P}[Z]=1\). Define
\[
\mathbb Q(A)=\int_A Z\,\mathrm{d}\mathbb P\qquad\text{for }A\in\mathcal G.
\]
\begin{examsubparts}
\item[\textnormal{a).}]
Show that \(\mathbb Q\) is a probability measure on \((\Omega,\mathcal G)\).
\item[\textnormal{b).}]
For a random variable~\(X\), we have
\[
\mathbb E_{\mathbb Q}[X\mid\mathcal F]=\frac{\mathbb E_{\mathbb P}[ZX\mid\mathcal F]}{\mathbb E_{\mathbb P}[Z\mid\mathcal F]}.
\]
\end{examsubparts}
\end{examproblems}
\end{document}
